%*************************************
% ภาคผนวก จ.
%*************************************
\appendixChapterTitle{จ}{ซอร์สโค้ดหลักของระบบ}
\vspace*{-0.8in}

% ── Local code listing configuration ──────────────
\definecolor{codebg}{HTML}{F5F5F5}
\definecolor{codegreen}{HTML}{2E7D32}
\definecolor{codeblue}{HTML}{1565C0}
\definecolor{codepurple}{HTML}{7B1FA2}
\definecolor{codegray}{HTML}{616161}
\definecolor{codestring}{HTML}{C62828}
\definecolor{codeborder}{HTML}{BDBDBD}

\lstdefinelanguage{JavaScript}{
  morekeywords={async,await,break,case,catch,class,const,continue,debugger,default,delete,
    do,else,export,extends,finally,for,function,if,import,in,instanceof,let,new,of,
    return,super,switch,this,throw,try,typeof,var,void,while,with,yield,from,as,
    true,false,null,undefined,static,get,set},
  sensitive=true,
  morecomment=[l]{//},
  morecomment=[s]{/*}{*/},
  morestring=[b]',
  morestring=[b]",
  morestring=[b]`,
}

\lstdefinelanguage{JSX}{
  morekeywords={async,await,break,case,catch,class,const,continue,debugger,default,delete,
    do,else,export,extends,finally,for,function,if,import,in,instanceof,let,new,of,
    return,super,switch,this,throw,try,typeof,var,void,while,with,yield,from,as,
    true,false,null,undefined,static,get,set},
  sensitive=true,
  morecomment=[l]{//},
  morecomment=[s]{/*}{*/},
  morestring=[b]',
  morestring=[b]",
  morestring=[b]`,
}

\lstset{
  basicstyle=\ttfamily\fontsize{9}{11}\selectfont,
  backgroundcolor=\color{codebg},
  keywordstyle=\color{codeblue}\bfseries,
  commentstyle=\color{codegreen}\itshape,
  stringstyle=\color{codestring},
  numberstyle=\tiny\color{codegray},
  numbers=left,
  numbersep=8pt,
  frame=single,
  rulecolor=\color{codeborder},
  tabsize=2,
  breaklines=true,
  breakatwhitespace=false,
  showstringspaces=false,
  showspaces=false,
  keepspaces=true,
  columns=flexible,
  xleftmargin=15pt,
  framexleftmargin=15pt,
  aboveskip=8pt,
  belowskip=8pt,
  captionpos=b,
  escapeinside={(*@}{@*)},
  literate={é}{{\'e}}1 {è}{{\`e}}1 {ê}{{\^e}}1,
}
% ── End local code listing configuration ──────────

\section{โครงสร้างระบบฝั่งเซิร์ฟเวอร์ (Backend)}

\hspace*{1.3em}
ในส่วนนี้แสดงซอร์สโค้ดหลักของระบบแบ่งตามโมดูล ประกอบด้วย Express Server, Authentication, LLM Service, SSH Service, NETCONF Service, Console Service, Agent Relay Service, API Routes, Vercel Serverless Handler และฝั่ง Frontend ที่ใช้ React

%% ─────────────────────────────────────
\subsection{Express Server (server.js)}
\hspace*{1.3em}
ไฟล์หลักของเซิร์ฟเวอร์ Express ที่ทำหน้าที่ตั้งค่า Middleware, Security, CORS, Routes และเชื่อมต่อ MongoDB Atlas รวมถึง Vercel Serverless Handler

\begin{lstlisting}[language=JavaScript, caption={server.js -- Express Server Setup}]
import express from 'express';
import cors from 'cors';
import helmet from 'helmet';
import rateLimit from 'express-rate-limit';
import mongoose from 'mongoose';
import session from 'express-session';
import { config } from './config/config.js';
import { mongooseOptions } from './lib/mongodb.js';
import { compressionMiddleware, responseTimeMiddleware }
  from './middleware/performance.js';
import devicesRouter from './routes/devices.js';
import configurationsRouter from './routes/configurations.js';
import consoleRouter from './routes/console.js';
import backupsRouter from './routes/backups.js';
import yangModelsRouter from './routes/yangModels.js';
import authRouter, { passport } from './routes/auth.js';
import agentRouter from './routes/agent.js';

const app = express();

// Security middleware
app.use(helmet({
  contentSecurityPolicy: false,
  crossOriginEmbedderPolicy: false
}));

// Rate limiting
const limiter = rateLimit({
  windowMs: 5 * 60 * 1000,
  max: 1000,
  message: { success: false,
    message: 'Too many requests, please try again later.'
  },
  skip: (req) => {
    return req.url.includes('/console')
      || req.url.includes('/health')
      || req.url.includes('/backups');
  }
});

if (config.server.nodeEnv === 'production') {
  app.use('/api/', limiter);
}

app.use(cors(config.cors));
app.use(compressionMiddleware);
app.use(responseTimeMiddleware);
app.use(express.json({ limit: '10mb' }));
app.use(express.urlencoded({ extended: true }));

// Session configuration (required for Passport)
app.use(session({
  secret: config.auth.sessionSecret,
  resave: false,
  saveUninitialized: false,
  cookie: {
    secure: config.server.nodeEnv === 'production',
    httpOnly: true,
    maxAge: 24 * 60 * 60 * 1000
  }
}));

app.use(passport.initialize());
app.use(passport.session());

// Health check endpoint
app.get('/api/health', async (req, res) => {
  try {
    if (mongoose.connection.readyState === 1) {
      await mongoose.connection.db.admin().ping();
    } else {
      throw new Error('MongoDB not connected');
    }
    res.json({
      success: true, message: 'Server is healthy',
      timestamp: new Date().toISOString(),
      version: '2.0.0',
      environment: config.server.nodeEnv,
      database: 'MongoDB Atlas'
    });
  } catch (error) {
    res.status(503).json({
      success: false, message: 'Server is unhealthy',
      error: error.message
    });
  }
});

// API routes
app.use('/api/auth', authRouter);
app.use('/api/devices', devicesRouter);
app.use('/api/configurations', configurationsRouter);
app.use('/api/console', consoleRouter);
app.use('/api/backups', backupsRouter);
app.use('/api/yang-models', yangModelsRouter);
app.use('/api/agent', agentRouter);

// Error handling middleware
app.use((err, req, res, next) => {
  console.error('Unhandled error:', err);
  res.status(err.status || 500).json({
    success: false,
    message: config.server.nodeEnv === 'production'
      ? 'Internal server error' : err.message
  });
});

// MongoDB connection
async function connectToMongoDB() {
  try {
    await mongoose.connect(
      config.database.mongodb_uri, mongooseOptions
    );
    console.log('Connected to MongoDB Atlas');
  } catch (error) {
    console.error('Failed to connect to MongoDB:', error);
    if (!process.env.VERCEL) process.exit(1);
    throw error;
  }
}

// Vercel serverless handler
export default async (req, res) => {
  if (process.env.VERCEL) await connectDB();
  return app(req, res);
};
\end{lstlisting}

%% ─────────────────────────────────────
\subsection{Authentication Middleware (auth.js)}
\hspace*{1.3em}
Middleware สำหรับตรวจสอบสิทธิ์ผู้ใช้ด้วย JWT Token พร้อมระบบ Cache เพื่อลดการอ่านฐานข้อมูลซ้ำ

\begin{lstlisting}[language=JavaScript, caption={middleware/auth.js -- JWT Authentication}]
import jwt from 'jsonwebtoken';
import { config } from '../config/config.js';
import User from '../models/User.js';

const userCache = new Map();
const USER_CACHE_TTL = 60 * 1000;

function getCachedUser(userId) {
  const entry = userCache.get(userId);
  if (entry && Date.now() - entry.timestamp < USER_CACHE_TTL)
    return entry.user;
  if (entry) userCache.delete(userId);
  return null;
}

function setCachedUser(userId, user) {
  userCache.set(userId, { user, timestamp: Date.now() });
  if (userCache.size > 500) {
    const oldest = userCache.keys().next().value;
    userCache.delete(oldest);
  }
}

export const authenticateToken = async (req, res, next) => {
  try {
    const authHeader = req.headers['authorization'];
    const token = authHeader && authHeader.split(' ')[1];
    if (!token) return res.status(401).json({
      success: false, message: 'Access token required'
    });

    const decoded = jwt.verify(token, config.auth.jwtSecret);
    let user = getCachedUser(decoded.userId);
    if (!user) {
      user = await User.findById(decoded.userId).lean();
      if (user) setCachedUser(decoded.userId, user);
    }
    if (!user) return res.status(401).json({
      success: false, message: 'User not found'
    });
    if (!user.isActive) return res.status(403).json({
      success: false, message: 'Account is deactivated'
    });

    req.user = user;
    req.userId = user._id;
    next();
  } catch (error) {
    if (error.name === 'TokenExpiredError')
      return res.status(401).json({
        success: false, message: 'Token expired',
        code: 'TOKEN_EXPIRED'
      });
    if (error.name === 'JsonWebTokenError')
      return res.status(401).json({
        success: false, message: 'Invalid token'
      });
    return res.status(500).json({
      success: false, message: 'Authentication error'
    });
  }
};

export const generateToken = (user) => {
  return jwt.sign(
    { userId: user._id, email: user.email, role: user.role },
    config.auth.jwtSecret,
    { expiresIn: config.auth.jwtExpiresIn || '7d' }
  );
};

export const generateRefreshToken = (user) => {
  return jwt.sign(
    { userId: user._id, type: 'refresh' },
    config.auth.jwtSecret,
    { expiresIn: '30d' }
  );
};
\end{lstlisting}

%% ─────────────────────────────────────
\subsection{LLM Service (llmService.js)}
\hspace*{1.3em}
บริการหลักสำหรับการเชื่อมต่อกับ Large Language Model (LLM) รองรับทั้ง OpenRouter (Cloud) และ Ollama (Local ผ่าน Agent) เพื่อสร้าง Configuration อัตโนมัติสำหรับอุปกรณ์ Cisco

\begin{lstlisting}[language=JavaScript, caption={services/llmService.js -- LLM Service Class}]
import axios from 'axios';
import agentRelay from './agentRelay.js';

export class LLMService {
  constructor() {
    this.provider = process.env.LLM_PROVIDER || 'openrouter';
    this.apiKey = process.env.OPENROUTER_API_KEY || '';
    this.model = process.env.OPENROUTER_MODEL || '';
    this.ollamaModel = process.env.OLLAMA_MODEL || 'llama3.2';
    this.apiUrl =
      'https://openrouter.ai/api/v1/chat/completions';
    this.timeout = 120000;

    this.client = axios.create({
      baseURL: 'https://openrouter.ai/api/v1',
      timeout: this.timeout,
      headers: {
        'Authorization': 'Bearer ' + this.apiKey,
        'Content-Type': 'application/json'
      }
    });

    this.defaultParams = {
      temperature: 0.1, max_tokens: 1000,
      top_p: 0.85, frequency_penalty: 0.3,
      presence_penalty: 0.2,
      stop: ['```', 'Here are', 'Sure', 'Note:']
    };

    this.cache = new Map();
    this.cacheTTL = 5 * 60 * 1000;
    this.maxCacheSize = 100;

    this.rateLimiter = {
      requests: [], maxRequests: 30,
      windowMs: 60 * 1000
    };

    this.knowledgeBase = this._initializeKnowledgeBase();
  }
\end{lstlisting}

\begin{lstlisting}[language=JavaScript, caption={llmService.js -- Chat Completion และ Ollama Relay}]
  async _callChatCompletion(messages, params = {}, userId = null) {
    if (this.provider === 'ollama') {
      return this._chatViaOllama(userId, messages, params);
    }
    const response = await this.client.post(
      '/chat/completions', {
      model: params.model || this.model,
      messages,
      temperature: params.temperature
        ?? this.defaultParams.temperature,
      max_tokens: params.max_tokens
        ?? this.defaultParams.max_tokens,
      top_p: params.top_p ?? this.defaultParams.top_p,
    });
    return response.data;
  }

  async _chatViaOllama(userId, messages, params = {}) {
    if (!userId)
      throw new Error('Ollama mode requires an active agent.');
    const isOnline = await agentRelay.isAgentOnline(userId);
    if (!isOnline)
      throw new Error('Agent not connected.');

    const result = await agentRelay.sendToAgent(
      userId, 'agent:ollama:chat', {
      messages,
      model: params.model || this.ollamaModel,
      temperature: params.temperature
        ?? this.defaultParams.temperature,
      max_tokens: params.max_tokens
        ?? this.defaultParams.max_tokens,
    }, 8000);

    if (result.pending) {
      return { pending: true, commandId: result.commandId };
    }
    return result;
  }
\end{lstlisting}

\begin{lstlisting}[language=JavaScript, caption={llmService.js -- generateConfiguration}]
  async generateConfiguration(
    prompt, deviceType, deviceContext = {},
    templateName = 'cisco_cli', useCache = true,
    options = {}
  ) {
    const startTime = Date.now();
    this.stats.totalRequests++;
    try {
      if (useCache) {
        const cacheKey = this._getCacheKey(
          prompt, deviceType, 'cli'
        );
        const cached = this._getFromCache(cacheKey);
        if (cached) return {
          ...cached, fromCache: true,
          executionTime: Date.now() - startTime
        };
      }
      this._checkRateLimit();

      let systemMessage, userMessage;
      const built = this.buildFromTemplate(
        templateName, prompt, deviceType, deviceContext
      );
      if (built.success) {
        systemMessage = built.system;
        userMessage = built.user;
      } else {
        systemMessage = this._buildSystemMessage(
          prompt, deviceType);
        userMessage = this._buildUserMessage(
          prompt, deviceType, deviceContext);
      }

      const messages = [
        { role: 'system', content: systemMessage },
        { role: 'user', content: userMessage }
      ];

      let llmResponse;
      let retryCount = 0;
      const maxRetries = this.provider === 'ollama' ? 0 : 2;
      while (retryCount <= maxRetries) {
        try {
          llmResponse = await this._callChatCompletion(
            messages, this.defaultParams,
            options.userId || null
          );
          break;
        } catch (apiError) {
          retryCount++;
          if (retryCount > maxRetries) throw apiError;
          await new Promise(r =>
            setTimeout(r, 1000 * retryCount));
        }
      }

      if (llmResponse.pending) {
        return {
          success: true, pending: true,
          commandId: llmResponse.commandId,
          message: 'Generation in progress via Ollama.'
        };
      }

      const rawConfig = llmResponse?.choices?.[0]
        ?.message?.content?.trim();
      if (!rawConfig) throw new Error('Empty response');

      const cleanConfig = this._cleanConfiguration(rawConfig);
      const ensuredConfig = this._ensureTrunkAllowed(
        cleanConfig, prompt);
      const deployConfig = this._createDeploymentVersion(
        ensuredConfig);
      const validation = this._validateConfiguration(
        ensuredConfig, deviceType);

      const result = {
        success: true,
        configuration: deployConfig,
        displayConfig: ensuredConfig,
        model: this.getActiveModel(),
        provider: this.provider,
        deviceType,
        executionTime: Date.now() - startTime,
        validation,
        confidenceScore: validation.score,
        fromCache: false
      };

      if (useCache) {
        const cacheKey = this._getCacheKey(
          prompt, deviceType, 'cli');
        this._setCache(cacheKey, result);
      }
      return result;
    } catch (error) {
      this.stats.errors++;
      return {
        success: false,
        error: 'Generation failed: ' + error.message,
        configuration: null,
        executionTime: Date.now() - startTime
      };
    }
  }
\end{lstlisting}

%% ─────────────────────────────────────
\subsection{SSH Service (sshService.js)}
\hspace*{1.3em}
บริการสำหรับเชื่อมต่อ SSH ไปยังอุปกรณ์เครือข่าย รองรับ Session Reuse, Persistent Sessions, การ Deploy Configuration ผ่าน Shell และ Privilege Escalation (Enable Mode)

\begin{lstlisting}[language=JavaScript, caption={services/sshService.js -- SSH Service Class}]
let Client;
try {
  const ssh2 = await import('ssh2');
  Client = ssh2.Client;
} catch (e) {
  console.warn('ssh2 not available (serverless mode)');
}

export class SSHService {
  constructor() {
    this.connections = new Map();
    this.persistentSessions = new Map();
    this.sessionTimeout = 600000;
    this.maxSessionsPerDevice = 3;
    this.commandQueue = new Map();
    this.sessionPool = new Map();

    this.cleanupInterval = setInterval(() => {
      this.cleanupExpiredSessions();
      this.optimizeSessionPool();
    }, 120000);

    this.healthCheckInterval = setInterval(() => {
      this.healthCheckSessions();
    }, 30000);
  }

  async connect(deviceConfig) {
    const { id, _id, ip_address, ssh_port,
            username, password } = deviceConfig;
    const deviceId = id || _id;
    this.disconnect(deviceId);

    return new Promise((resolve, reject) => {
      const conn = new Client();
      const timeout = setTimeout(() => {
        conn.end();
        reject(new Error('SSH connection timeout'));
      }, 10000);

      conn.on('ready', () => {
        clearTimeout(timeout);
        conn._isReady = true;
        this.connections.set(deviceId, {
          connection: conn,
          createdAt: Date.now(),
          lastUsed: Date.now(),
          isReady: true
        });
        resolve(conn);
      });

      conn.on('error', (err) => {
        clearTimeout(timeout);
        conn._isReady = false;
        this.connections.delete(deviceId);
        reject(err);
      });

      conn.connect({
        host: ip_address,
        port: ssh_port || 22,
        username, password,
        readyTimeout: 15000,
        authTimeout: 10000,
        tryKeyboard: true,
        keepaliveInterval: 15000,
        keepaliveCountMax: 5,
        algorithms: {
          kex: ['diffie-hellman-group14-sha1',
                'ecdh-sha2-nistp256'],
          cipher: ['aes128-ctr', 'aes256-ctr',
                   'aes128-cbc', '3des-cbc'],
        }
      });
    });
  }
\end{lstlisting}

\begin{lstlisting}[language=JavaScript, caption={sshService.js -- sendConfigCommands (Deploy Configuration)}]
  async sendConfigCommands(deviceConfig, commands,
    options = {}
  ) {
    const { tolerateErrors = false,
            useSession = true } = options;
    let conn, sessionReused = false;
    const deviceId = deviceConfig.id || deviceConfig._id;
    const sessionKey = deviceId + '_persistent';

    // Try reuse existing persistent session
    if (useSession) {
      const existing =
        this.persistentSessions.get(sessionKey);
      if (existing && this.isSessionValid(existing)) {
        conn = existing.connection;
        existing.lastUsed = Date.now();
        existing.useCount++;
        sessionReused = true;
      } else {
        conn = await this.connect(deviceConfig);
        this.persistentSessions.set(sessionKey, {
          connection: conn, deviceId, deviceConfig,
          createdAt: Date.now(),
          lastUsed: Date.now(), useCount: 1
        });
      }
    } else {
      conn = await this.connect(deviceConfig);
    }

    // Create shell
    let stream;
    try {
      stream = await new Promise((resolve, reject) => {
        conn.shell({ pty: true },
          (err, s) => err ? reject(err) : resolve(s));
      });
    } catch (shellError) {
      this.persistentSessions.delete(sessionKey);
      conn = await this.connect(deviceConfig);
      stream = await new Promise((resolve, reject) => {
        conn.shell({ pty: true },
          (err, s) => err ? reject(err) : resolve(s));
      });
    }

    return new Promise((resolve, reject) => {
      let output = '', currentStep = 0;
      let errorCount = 0, hasErrors = false;
      const configCommands = commands.split('\n')
        .map(c => c.trim())
        .filter(c => c
          && !c.startsWith('configure terminal')
          && !c.startsWith('end'));
      const allCommands = [
        'enable', 'configure terminal',
        ...configCommands, 'end', 'write memory'
      ];

      const timeout = setTimeout(() => {
        stream.end();
        reject(new Error('Deployment timeout'));
      }, 30000);

      const sendNextCommand = () => {
        if (currentStep >= allCommands.length) {
          clearTimeout(timeout);
          setTimeout(() => {
            stream.end();
            resolve({
              success: true, output: output.trim(),
              commandsExecuted: allCommands,
              errorCount, sessionReused
            });
          }, 2000);
          return;
        }
        stream.write(
          allCommands[currentStep] + '\r\n');
        currentStep++;
      };

      stream.on('data', (data) => {
        const chunk = data.toString();
        output += chunk;
        if (chunk.includes('#')
          || chunk.includes('Password:')
          || chunk.includes('(config)#')) {
          setTimeout(sendNextCommand, 500);
        }
        if (chunk.toLowerCase().includes('password:'))
          stream.write(
            deviceConfig.password + '\r\n');
        if (chunk.includes('% Invalid')
          || chunk.includes('% Error')) {
          errorCount++; hasErrors = true;
        }
      });
    });
  }
\end{lstlisting}

%% ─────────────────────────────────────
\subsection{NETCONF Service (netconfService.js)}
\hspace*{1.3em}
บริการสำหรับการเชื่อมต่อ NETCONF (RFC 6241) ผ่าน SSH Subsystem รองรับทั้ง Cisco NX-OS (Nexus) และ IOS-XE พร้อม YANG Model Operations

\begin{lstlisting}[language=JavaScript, caption={services/netconfService.js -- NETCONF Service Class}]
let Client;
try {
  const ssh2 = await import('ssh2');
  Client = ssh2.Client;
} catch (e) {
  console.warn('ssh2 not available (serverless mode)');
}

export class NetconfService {
  constructor() {
    this.connections = new Map();
    this.defaultPort = 830;
    this.sessionTimeout = 300000;

    this.NAMESPACES = {
      'nxos':
        'http://cisco.com/ns/yang/cisco-nx-os-device',
      'ios-xe':
        'http://cisco.com/ns/yang/Cisco-IOS-XE-native',
    };

    this.ROOT_ELEMENTS = {
      'nxos': 'System', 'nexus': 'System',
      'ios-xe': 'native', 'ios': 'native'
    };

    this.NETCONF_HELLO = '<?xml version="1.0"?>\n'
      + '<hello xmlns="urn:ietf:params:xml:ns:'
      + 'netconf:base:1.0">\n  <capabilities>\n'
      + '    <capability>urn:ietf:params:netconf:'
      + 'base:1.0</capability>\n'
      + '    <capability>urn:ietf:params:netconf:'
      + 'capability:writable-running:1.0'
      + '</capability>\n'
      + '    <capability>urn:ietf:params:netconf:'
      + 'capability:candidate:1.0</capability>\n'
      + '  </capabilities>\n</hello>]]>]]>';

    this.MESSAGE_DELIMITER = ']]>]]>';
  }
\end{lstlisting}

\begin{lstlisting}[language=JavaScript, caption={netconfService.js -- connect}]
  async connect(deviceConfig) {
    if (!Client) throw new Error(
      'NETCONF not available. Use agent relay.');
    const { id, _id, ip_address, username,
            password, netconf_port } = deviceConfig;
    const deviceId = id || _id;
    const port = netconf_port || this.defaultPort;
    this.disconnect(deviceId);

    return new Promise((resolve, reject) => {
      const conn = new Client();
      const timeout = setTimeout(() => {
        conn.end();
        reject(new Error('NETCONF connection timeout'));
      }, 30000);

      conn.on('ready', () => {
        clearTimeout(timeout);
        conn.subsys('netconf', (err, stream) => {
          if (err) { conn.end(); reject(err); return; }
          let serverHello = '';
          let helloReceived = false;

          const helloHandler = (data) => {
            serverHello += data.toString();
            if (serverHello.includes(
              this.MESSAGE_DELIMITER)
              && !helloReceived
            ) {
              helloReceived = true;
              stream.removeListener(
                'data', helloHandler);
              const capabilities =
                this.parseCapabilities(serverHello);
              stream.write(this.NETCONF_HELLO,
                'utf8', (err) => {
                if (err) { reject(err); return; }
                this.connections.set(deviceId, {
                  connection: conn, stream,
                  capabilities,
                  createdAt: Date.now(),
                  lastUsed: Date.now()
                });
                resolve({
                  success: true, deviceId,
                  capabilities,
                  message: 'NETCONF session established'
                });
              });
            }
          };
          stream.on('data', helloHandler);
        });
      });

      conn.connect({
        host: ip_address, port,
        username, password,
        readyTimeout: 30000,
        algorithms: {
          kex: ['diffie-hellman-group14-sha256',
                'ecdh-sha2-nistp256'],
          cipher: ['aes256-ctr', 'aes128-ctr']
        }
      });
    });
  }
\end{lstlisting}

\begin{lstlisting}[language=JavaScript, caption={netconfService.js -- sendRpc และ editConfig}]
  async sendRpc(deviceId, rpcContent,
    messageId = null, timeoutMs = 60000
  ) {
    const session = this.connections.get(deviceId);
    if (!session)
      throw new Error('No active NETCONF session');
    session.lastUsed = Date.now();
    const msgId = messageId || 'msg-' + Date.now();

    const rpcMessage =
      '<?xml version="1.0" encoding="UTF-8"?>\n'
      + '<rpc message-id="' + msgId
      + '" xmlns="urn:ietf:params:xml:ns:'
      + 'netconf:base:1.0">\n'
      + rpcContent
      + '\n</rpc>' + this.MESSAGE_DELIMITER;

    return new Promise((resolve, reject) => {
      let response = '';
      const timeout = setTimeout(() => {
        session.stream.removeListener(
          'data', dataHandler);
        reject(new Error('NETCONF timeout'));
      }, timeoutMs);

      const dataHandler = (data) => {
        response += data.toString();
        if (response.includes(
          this.MESSAGE_DELIMITER)) {
          clearTimeout(timeout);
          session.stream.removeListener(
            'data', dataHandler);
          const clean = response
            .replace(this.MESSAGE_DELIMITER, '')
            .trim();
          if (clean.includes('<rpc-error>')) {
            reject(new Error('NETCONF error: '
              + this.parseRpcError(clean)));
          } else {
            resolve({
              success: true,
              messageId: msgId,
              response: clean
            });
          }
        }
      };
      session.stream.on('data', dataHandler);
      session.stream.write(rpcMessage, 'utf8');
    });
  }

  async editConfig(deviceId, configXml,
    target = 'running',
    defaultOperation = 'merge'
  ) {
    const rpc = '  <edit-config>'
      + '<target><' + target + '/></target>'
      + '<default-operation>' + defaultOperation
      + '</default-operation>'
      + '<config>\n' + configXml
      + '\n    </config></edit-config>';
    return await this.sendRpc(deviceId, rpc, null, 90000);
  }

  async getRunningConfig(deviceId, filter = null) {
    let filterXml = filter
      ? '\n  <filter type="subtree">'
        + filter + '</filter>'
      : '';
    return await this.sendRpc(deviceId,
      '  <get-config><source><running/></source>'
      + filterXml + '</get-config>', null, 180000);
  }
\end{lstlisting}

%% ─────────────────────────────────────
\subsection{Console Service (consoleService.js)}
\hspace*{1.3em}
บริการสำหรับเชื่อมต่อผ่าน Serial Console (RS-232/USB) เพื่อทำ Initial Configuration ของอุปกรณ์เครือข่ายที่ยังไม่มี IP Address

\begin{lstlisting}[language=JavaScript, caption={services/consoleService.js -- Console Service}]
let SerialPort, ReadlineParser;
try {
  const sp = await import('serialport');
  SerialPort = sp.SerialPort;
  const rp = await import(
    '@serialport/parser-readline');
  ReadlineParser = rp.ReadlineParser;
} catch (e) {
  console.warn('serialport not available');
}

export class ConsoleService {
  constructor() {
    this.connections = new Map();
    this.serialPortAvailable = !!SerialPort;
    this.defaultSettings = {
      baudRate: 9600, dataBits: 8,
      parity: 'none', stopBits: 1,
      flowControl: false
    };
  }

  async getAvailablePorts() {
    const ports = await SerialPort.list();
    const formattedPorts = ports.map(port => {
      let displayName = port.path;
      if (port.manufacturer) {
        if (port.manufacturer.toLowerCase()
          .includes('ftdi'))
          displayName = port.path
            + ' - FTDI USB Serial';
        else
          displayName = port.path
            + ' - ' + port.manufacturer;
      }
      return {
        path: port.path,
        manufacturer: port.manufacturer || 'Unknown',
        vendorId: port.vendorId || 'N/A',
        friendlyName: displayName,
        isUSB: port.path.includes('USB')
      };
    });
    return {
      success: true, ports: formattedPorts,
      count: formattedPorts.length
    };
  }

  async connectConsole(connectionConfig) {
    const { deviceId, portPath, baudRate } =
      connectionConfig;
    if (this.connections.has(deviceId))
      await this.disconnectConsole(deviceId);

    const port = new SerialPort({
      path: portPath,
      baudRate: baudRate || 9600,
      dataBits: 8, parity: 'none', stopBits: 1,
      autoOpen: false
    });
    const parser = port.pipe(
      new ReadlineParser({ delimiter: '\r\n' }));

    return new Promise((resolve, reject) => {
      port.open((err) => {
        if (err) { reject(err); return; }
        this.connections.set(deviceId, {
          port, parser, portPath,
          connectedAt: new Date()
        });
        resolve({
          success: true,
          message: 'Console connected via ' + portPath
        });
      });
    });
  }

  async sendInitialConfig(deviceId, configCommands) {
    const connection = this.connections.get(deviceId);
    if (!connection)
      throw new Error('No active console connection');
    let configResults = [];
    const commands = configCommands.split('\n')
      .map(l => l.trim())
      .filter(l => l && !l.startsWith('#'));

    for (let i = 0; i < commands.length; i++) {
      try {
        const result = await this.sendConsoleCommand(
          deviceId, commands[i], true);
        configResults.push({
          command: commands[i], success: true,
          output: result.output, sequence: i + 1
        });
        await new Promise(r => setTimeout(r, 500));
      } catch (cmdError) {
        configResults.push({
          command: commands[i], success: false,
          error: cmdError.message, sequence: i + 1
        });
      }
    }

    const successCount =
      configResults.filter(r => r.success).length;
    return {
      success: true,
      message: 'Initial configuration completed: '
        + successCount + '/' + configResults.length,
      results: configResults
    };
  }
}

export default new ConsoleService();
\end{lstlisting}

%% ─────────────────────────────────────
\subsection{Agent Relay Service (agentRelay.js)}
\hspace*{1.3em}
บริการ Relay สำหรับสื่อสารระหว่าง Backend (Vercel Serverless) กับ Desktop Agent ผ่าน MongoDB Command Queue โดยใช้ HTTP Polling เป็นกลไกการสื่อสาร

\begin{lstlisting}[language=JavaScript, caption={services/agentRelay.js -- Agent Relay (HTTP Polling)}]
import AgentCommand from '../models/AgentCommand.js';
import AgentHeartbeat from '../models/AgentHeartbeat.js';
import Notification from '../models/Notification.js';

class AgentRelay {
  constructor() {}

  initialize() {
    console.log('Agent relay initialized');
  }

  async sendToAgent(userId, event, data,
    timeoutMs = 8000
  ) {
    const online = await this.isAgentOnline(userId);
    if (!online) throw new Error('No agent connected.');

    const command = await AgentCommand.create({
      userId, event, data, status: 'pending'
    });
    const maxWait = Math.min(timeoutMs, 55000);
    const pollInterval = 300;
    const startTime = Date.now();

    while (Date.now() - startTime < maxWait) {
      const updated = await AgentCommand
        .findById(command._id);
      if (updated.status === 'completed')
        return {
          success: true, ...(updated.result || {})
        };
      if (updated.status === 'failed')
        throw new Error(
          updated.error || 'Command failed');
      await new Promise(r =>
        setTimeout(r, pollInterval));
    }

    return {
      success: false, pending: true,
      commandId: command._id.toString(),
      message: 'Command sent. Poll for result.'
    };
  }

  async isAgentOnline(userId) {
    const agent = await AgentHeartbeat
      .getOnlineAgent(userId.toString());
    return !!agent;
  }

  async getAgentInfo(userId) {
    return await AgentHeartbeat
      .getOnlineAgent(userId.toString());
  }
}

const agentRelay = new AgentRelay();
export default agentRelay;
\end{lstlisting}

%% ─────────────────────────────────────
\subsection{Device Routes (routes/devices.js)}
\hspace*{1.3em}
Express Routes สำหรับจัดการอุปกรณ์เครือข่าย (CRUD) พร้อม Joi Validation, Duplicate Check และ Cache Invalidation

\begin{lstlisting}[language=JavaScript, caption={routes/devices.js -- Device Validation Schema และ CRUD}]
import express from 'express';
import Joi from 'joi';
import Device from '../models/Device.js';
import sshService from '../services/sshService.js';
import netconfService
  from '../services/netconfService.js';
import { authenticateToken }
  from '../middleware/auth.js';
import { getOrSetCache, invalidateCache, CacheKeys }
  from '../lib/cache.js';

const router = express.Router();
router.use(authenticateToken);

const deviceSchema = Joi.object({
  name: Joi.string().required().max(255),
  type: Joi.string().valid(
    'router', 'switch', 'nexus', 'ios-xe'
  ).required(),
  layer: Joi.string().valid('layer-2', 'layer-3')
    .optional().allow('', null),
  ip_address: Joi.string().ip().required(),
  ssh_port: Joi.number().integer()
    .min(1).max(65535).default(22),
  netconf_port: Joi.number().integer()
    .min(1).max(65535).default(830),
  netconf_enabled: Joi.boolean().default(false),
  username: Joi.string().required().max(255),
  password: Joi.string().required().max(255),
  description: Joi.string().allow('').max(1000),
  location: Joi.string().allow('').max(255),
  model: Joi.string().allow('').max(255),
  status: Joi.string().valid(
    'active', 'inactive', 'maintenance', 'error'
  ).default('active')
}).options({ stripUnknown: true });

// GET /api/devices
router.get('/', async (req, res) => {
  try {
    const { status, type, limit = 50,
            offset = 0 } = req.query;
    const filter = { userId: req.userId };
    if (status) filter.status = status;
    if (type) filter.type = type;

    const devices = await Device.aggregate([
      { $match: filter },
      { $sort: { createdAt: -1 } },
      { $skip: parseInt(offset) },
      { $limit: parseInt(limit) },
      { $lookup: {
          from: 'configurationhistories',
          localField: '_id',
          foreignField: 'device_id',
          as: 'config_history'
      }},
      { $addFields: {
          id: '$_id',
          total_configs: {
            $size: '$config_history' },
          deployed_configs: {
            $size: { $filter: {
              input: '$config_history',
              as: 'c',
              cond: { $eq: [
                '$$c.status', 'deployed'] }
            }}
          }
      }}
    ]);
    res.json({ success: true, devices });
  } catch (error) {
    res.status(500).json({
      success: false,
      message: 'Failed to fetch devices'
    });
  }
});

// PUT /api/devices/:id
router.put('/:id', async (req, res) => {
  try {
    const { id } = req.params;
    const { error, value } =
      deviceUpdateSchema.validate(req.body);
    if (error) return res.status(400).json({
      success: false, message: 'Validation error',
      details: error.details
    });

    if (value.password === '') delete value.password;
    if (value.type && value.type !== 'switch')
      delete value.layer;

    const device = await Device.findOneAndUpdate(
      { _id: id, userId: req.userId },
      { ...value, updatedAt: new Date() },
      { new: true }
    );
    if (!device) return res.status(404).json({
      success: false, message: 'Device not found'
    });

    invalidateCache(CacheKeys.devices(req.userId));
    res.json({
      success: true, device,
      message: 'Device updated successfully'
    });
  } catch (error) {
    res.status(500).json({
      success: false,
      message: 'Failed to update device'
    });
  }
});
\end{lstlisting}

%% ─────────────────────────────────────
\subsection{Configuration Deploy via Agent (routes/configurations.js)}
\hspace*{1.3em}
ฟังก์ชัน Deploy Configuration ผ่าน Agent Relay โดยตรวจสอบว่ามี Agent ออนไลน์หรือไม่ ถ้ามีจะส่งคำสั่งผ่าน Agent แทนการเชื่อมต่อ SSH โดยตรง

\begin{lstlisting}[language=JavaScript, caption={routes/configurations.js -- Deploy via Agent}]
import express from 'express';
import Device from '../models/Device.js';
import ConfigurationHistory
  from '../models/ConfigurationHistory.js';
import llmService from '../services/llmService.js';
import sshService from '../services/sshService.js';
import agentRelay from '../services/agentRelay.js';
import { authenticateToken }
  from '../middleware/auth.js';

const router = express.Router();

async function deployViaAgent(
  userId, device, configCommands
) {
  const agentOnline =
    await agentRelay.isAgentOnline(userId);

  if (agentOnline) {
    const result = await agentRelay.sendToAgent(
      userId, 'agent:ssh:deploy-config', {
        deviceId: device._id.toString(),
        host: device.ip_address,
        port: device.ssh_port || 22,
        username: device.username,
        password: device.password,
        commands: configCommands,
        enablePassword: device.enable_password
      }, 25000);

    if (result.pending) return {
      success: true, pending: true,
      output: 'Configuration sent to agent...',
      commandId: result.commandId
    };
    if (!result.success)
      throw new Error(
        result.error || 'Agent deploy failed');
    return result;
  }

  // Fallback to direct SSH
  return await sshService.sendConfigCommands(
    device, configCommands);
}

router.use(authenticateToken);
\end{lstlisting}

%% ─────────────────────────────────────
\subsection{Vercel Serverless Handler (api/handler.js)}
\hspace*{1.3em}
Entry point สำหรับ Vercel Serverless Functions ทำหน้าที่ Lazy-load Express Application และ Proxy Request

\begin{lstlisting}[language=JavaScript, caption={api/handler.js -- Vercel Serverless Handler}]
let app;

async function getApp() {
  if (!app) {
    const mod = await import('../backend/server.js');
    app = mod.default;
  }
  return app;
}

export default async function handler(req, res) {
  try {
    const appHandler = await getApp();
    if (!req.url.startsWith('/api')) {
      req.url = '/api' + req.url;
    }
    return await appHandler(req, res);
  } catch (err) {
    console.error('Request error:', err.message);
    res.statusCode = 500;
    res.end(JSON.stringify({
      error: 'Internal server error'
    }));
  }
}

export const config = { maxDuration: 60 };
\end{lstlisting}

%% ─────────────────────────────────────
\section{โครงสร้างระบบฝั่งไคลเอนต์ (Frontend)}

\hspace*{1.3em}
ระบบ Frontend พัฒนาด้วย React + Vite โดยใช้ React Router สำหรับ Routing, Lazy Loading สำหรับ Code Splitting และ Axios สำหรับ HTTP Client

%% ─────────────────────────────────────
\subsection{React Entry Point (main.jsx)}

\begin{lstlisting}[language=JSX, caption={frontend/src/main.jsx -- React Application Entry}]
import { StrictMode } from 'react'
import { createRoot } from 'react-dom/client'
import { SpeedInsights }
  from '@vercel/speed-insights/react'
import './index.css'
import App from './App.jsx'
import { ThemeProvider }
  from './context/ThemeContext.jsx'

createRoot(document.getElementById('root')).render(
  <StrictMode>
    <ThemeProvider>
      <App />
      <SpeedInsights />
    </ThemeProvider>
  </StrictMode>,
)
\end{lstlisting}

%% ─────────────────────────────────────
\subsection{Application Router (App.jsx)}

\begin{lstlisting}[language=JSX, caption={frontend/src/App.jsx -- React Router และ Lazy Loading}]
import { HashRouter as Router, Routes, Route,
  Navigate } from 'react-router-dom';
import { useState, useEffect, lazy, Suspense }
  from 'react';
import { Toaster } from 'react-hot-toast';
import axios from 'axios';
import Sidebar from './components/Sidebar';
import ErrorBoundary
  from './components/ErrorBoundary';
import ProtectedRoute
  from './components/ProtectedRoute';
import PageTransition
  from './components/PageTransition';
import { AuthProvider, useAuth }
  from './context/AuthContext';
import Dashboard from './pages/Dashboard';
import Login from './pages/Login';

// Lazy-loaded pages (Code Splitting)
const Devices = lazy(() =>
  import('./pages/Devices'));
const Configurations = lazy(() =>
  import('./pages/Configurations'));
const ConfigurationHistory = lazy(() =>
  import('./pages/ConfigurationHistory'));
const ConsoleConfiguration = lazy(() =>
  import('./pages/ConsoleConfiguration'));
const BackupManagement = lazy(() =>
  import('./pages/BackupManagement'));
const Diagrams = lazy(() =>
  import('./pages/Diagrams'));
const AgentSettings = lazy(() =>
  import('./pages/AgentSettings'));

const API_BASE_URL = import.meta.env.VITE_API_URL
  || '/api';
axios.defaults.baseURL = API_BASE_URL;

function AuthenticatedLayout({ children,
  connectionStatus, sidebarCollapsed,
  onToggleSidebar
}) {
  return (
    <div className="flex min-h-screen">
      <Sidebar
        connectionStatus={connectionStatus}
        isCollapsed={sidebarCollapsed}
        onToggleCollapse={onToggleSidebar}
      />
      <div className="flex-1">
        <main className="pt-20 pb-8 px-4">
          <div className="max-w-7xl mx-auto">
            <PageTransition>
              {children}
            </PageTransition>
          </div>
        </main>
      </div>
    </div>
  );
}
\end{lstlisting}

%% ─────────────────────────────────────
\section{โครงสร้าง Desktop Agent (Electron)}

\hspace*{1.3em}
Desktop Agent เป็นแอปพลิเคชัน Electron ที่ทำหน้าที่เป็นตัวกลางระหว่าง Backend (Vercel) กับอุปกรณ์เครือข่ายในเครือข่ายภายใน โดยใช้ HTTP Polling เป็นกลไกการสื่อสาร รองรับ SSH, NETCONF, Serial Console และ Ollama LLM

%% ─────────────────────────────────────
\subsection{Agent Configuration และ HTTP Polling Client}

\begin{lstlisting}[language=JavaScript, caption={agent-gui/src/main.js -- AgentConfig และ HttpPollingClient}]
const { app, BrowserWindow, Tray, Menu,
  ipcMain } = require('electron');
const path = require('path');
const os = require('os');
const crypto = require('crypto');
const { EventEmitter } = require('events');
const fs = require('fs');

class AgentConfig {
  constructor() {
    this.configDir = path.join(
      app.getPath('userData'));
    this.configPath = path.join(
      this.configDir, 'config.json');
    this.data = this._load();
    if (!this.data.agentId) {
      this.data.agentId = crypto.randomBytes(16)
        .toString('hex');
      this._save();
    }
  }

  _load() {
    try {
      if (fs.existsSync(this.configPath))
        return JSON.parse(
          fs.readFileSync(
            this.configPath, 'utf8'));
    } catch (e) {}
    return {};
  }

  _save() {
    try {
      if (!fs.existsSync(this.configDir))
        fs.mkdirSync(
          this.configDir, { recursive: true });
      fs.writeFileSync(this.configPath,
        JSON.stringify(this.data, null, 2));
    } catch (e) {}
  }

  get(key) { return this.data[key] || ''; }
  set(key, value) {
    this.data[key] = value; this._save();
  }
}

class HttpPollingClient extends EventEmitter {
  constructor({ serverUrl, agentToken,
    agentName, version, handlers
  }) {
    super();
    this.serverUrl = serverUrl.replace(/\/$/, '');
    this.agentToken = agentToken;
    this.agentName = agentName;
    this.version = version;
    this.handlers = handlers;
    this.connected = false;
    this.pollInterval = null;
    this.heartbeatInterval = null;
    this.cachedPorts = [];
  }

  async connect() {
    await this._sendHeartbeat();
    this.connected = true;
    this.emit('connected');

    // Heartbeat every 5 seconds
    this.heartbeatInterval = setInterval(() => {
      this._sendHeartbeat()
        .catch(err => this.emit('error', err));
    }, 5000);

    // Poll commands every 2 seconds
    this.pollInterval = setInterval(() => {
      this._pollCommands().catch(() => {});
    }, 2000);

    // Refresh serial ports every 30 seconds
    this.portRefreshInterval =
      setInterval(async () => {
      try {
        const result =
          await this.handlers.console.listPorts();
        if (result && result.ports)
          this.cachedPorts = result.ports;
      } catch (e) {}
    }, 30000);
  }

  async _sendHeartbeat() {
    const response = await this._fetch(
      '/api/agent/poll/heartbeat', {
      method: 'POST',
      body: JSON.stringify({
        agentName: this.agentName,
        agentVersion: this.version,
        platform: os.platform(),
        hostname: os.hostname(),
        capabilities: {
          serialPorts: this.cachedPorts
        }
      })
    });
    if (!response.ok) throw new Error(
      'Heartbeat failed: ' + response.status);
  }

  async _pollCommands() {
    const response = await this._fetch(
      '/api/agent/poll/commands',
      { method: 'GET' });
    if (!response.ok) return;
    const { commands } = await response.json();
    for (const cmd of commands) {
      this.emit('command', {
        type: cmd.event,
        deviceId: cmd.data?.deviceId
      });
      this._handleCommand(cmd);
    }
  }

  async _handleCommand(cmd) {
    const { commandId, event, data } = cmd;
    try {
      let result;
      switch (event) {
        case 'agent:ssh:connect':
          result = await this.handlers
            .ssh.connect(data);
          break;
        case 'agent:ssh:deploy-config': {
          const { deviceId, host, port,
            username, password, commands } = data;
          await this.handlers.ssh.connect({
            deviceId, host,
            port: port || 22,
            username, password
          });
          result = await this.handlers
            .ssh.sendConfig(
              deviceId, commands);
          this.handlers.ssh.disconnect(deviceId);
          break;
        }
        case 'agent:netconf:connect':
          result = await this.handlers
            .netconf.connect(data);
          break;
        case 'agent:ollama:chat':
          result = await this.handlers
            .ollama.chat(data);
          break;
        case 'list-ports':
          result = await this.handlers
            .console.listPorts();
          break;
      }
      await this._reportResult(commandId, result);
    } catch (error) {
      await this._reportResult(
        commandId, null, error.message);
    }
  }
}
\end{lstlisting}
